%----------------------------------------------------------------------------------------
%	PACKAGES AND THEMES
%----------------------------------------------------------------------------------------

\PassOptionsToPackage{table}{xcolor}
\documentclass[aspectratio=169,xcolor=dvipsnames,svgnames,x11names,fleqn]{beamer}
% \documentclass[aspectratio=169,xcolor=dvipsnames,fleqn]{beamer}

\usetheme{RedVelvet}
\usefonttheme[onlymath]{serif}
\usepackage{xspace}
\usepackage{amsmath}
\usepackage{amssymb}
\usepackage{amsfonts}
\usepackage{color}
\usepackage{physics}
% \usepackage{mathbb}
\usepackage{rahul_math}
\usepackage{bigints}
\usepackage{hyperref}
\hypersetup{
colorlinks,
citecolor=Violet,
linkcolor=Black}
\usepackage{graphicx} % Allows including images
\usepackage{booktabs} % Allows the use of \toprule, \midrule and \bottomrule in tables
\usepackage{tikz,pgfplots}
\usepackage{subfigure}
\usetikzlibrary{arrows}
\usepackage{minted}
\definecolor{LightGray}{gray}{0.9}
\definecolor{cream}{rgb}{0.92, 0.9, 0.55}
\definecolor{lightblue}{rgb}{0.68, 0.85, 0.9}


\usepackage{xcolor-material}
\usetikzlibrary{fit}
\tikzset{%
apple/.pic={
    \fill [MaterialBrown] (-1/8,0)  arc (180:120:1 and 3/2) coordinate [pos=3/5] (@)-- ++(1/6,-1/7)  arc (120:180:5/4 and 3/2) -- cycle;
    \fill [MaterialLightGreen500] (0,-9/10)  .. controls ++(180:1/8) and ++(  0:1/4) .. (-1/3,  -1) .. controls ++(180:1/3) and ++(270:1/2) .. (  -1,   0) .. controls ++( 90:1/3) and ++(180:1/3) .. (-1/2, 3/4) .. controls ++(  0:1/8) and ++(135:1/8) .. (   0, 4/7)
}
}

\newcommand{\leftdoublequote}{\textcolor{blue}{\scalebox{3}{``}}}

\newcommand{\rightdoublequote}{\textcolor{blue}{\scalebox{3}{''}}}


\usepackage{textcomp}
\usepackage{fontawesome}

\usepackage{overpic}

%----------------------------------------------------------------------------------------
%	TITLE PAGE
%----------------------------------------------------------------------------------------

\usepackage{tikz-qtree,tikz-qtree-compat}
\usetikzlibrary{calc}


\title[CPE 381: Signals and Systems]{CPE 381: Fundamentals of Signals and Systems for Computer Engineers} % The short title appears at the bottom of every slide, the full title is only on the title page
\subtitle{03 Continuous-Time Signals and Systems}

\author[Rahul Bhadani] {{\Large \textbf{Rahul Bhadani}}}

\institute[UAH] % Your institution as it will appear on the bottom of every slide, maybe shorthand to save space
{
    Electrical \& Computer Engineering,  The University of Alabama in Huntsville
}
\date

% \titlegraphic{
%    \includegraphics[width=0.4\linewidth]{figures/UAH_primary.png}
% }
    
\begin{document}

%-------------------------------------------------
\begin{frame}
    \titlepage
\end{frame}

%-------------------------------------------------
\begin{frame}{Outline}
    \tableofcontents
\end{frame}

%%%%%%%%%%%%%%%%%%%%%%%%%%%%%%%%%%%%%%%%%%%%%%%%%%%%%
\section{Motivation}

%-----------------------------------------

\begin{frame}{}
    \begin{center}
    \Huge \bf \color{DarkBlue}
    \faFire
    
    Motivation
\end{center}
\end{frame}


\begin{frame}{}
    \begin{center}
    \Huge \bf 
    
    Signals and Systems is `Grandfather' of Data Science for Electrical and Computer Engineers

\end{center}
\end{frame}

\section{Continuous-Time Signals}

\begin{frame}
    \sectionpage
\end{frame}

\begin{frame}{Classification of Signals}

We care about the following properties when dealing with signals:

\begin{itemize}
    \item Predictability:  Random or Deterministic
    \item Variations of time and amplitude: continuous, discrete (time or x-axis) / quantized (amplitude or y-axis)
    \item Periodic/Aperiodic
    \item Finite energy/finite power;  Infinite energy/Infinite power

\end{itemize}

\end{frame}

%%%%%%%%%%%%%%%%%%%%%%%%%%%%%%%%%%%%%%%%%%%%%%%%%%%%%
\subsection{Operation on Signals}

%-----------------------------------------

\begin{frame}
    \subsectionpage
\end{frame}

\begin{frame}{Basic Mathematical Operations}
\begin{itemize}
    \item Addition: $x(t) + y(t)$
\item Subtraction: $x(t) - y(t)$
\item Constant multiplication: $kx(t)$ where $k$ is a constant
\end{itemize}
\end{frame}

\begin{frame}{Time-shift}

\begin{itemize}
    \item $x(t - \tau) \to $ Signal is delayed
\item $x(t + \tau) \to $ Signal is advanced
\end{itemize}
    \includegraphics[width=0.8\linewidth,trim=0 0 0 0cm,clip]{figures/Signal_Delay.png}
    
\end{frame}

\begin{frame}{Time Reflection}
\begin{columns}
    \column{0.3\linewidth}
    \begin{itemize}
    \item $x(t) \to x(-t)$ : take  mirror image\textbf{ along the y-axis}
    \end{itemize}
        \column{0.7\linewidth}
        Note: The book doesn’t specify whether to take the mirror image along the y-axis or not and it is confusing because the signal used in example 1.3.1 is symmetric with respect to both the x and y axes.
        \includegraphics[width=0.9\linewidth,trim=0 0 0 0cm,clip]{figures/Time_Reflection.png}
\end{columns}
\end{frame}


\begin{frame}{Signal Stretching along $y$-axis}

    \begin{itemize}
    \item $f(x) \to a f(x); \quad a > 1$ : Shrink the graph of $f(x)$ `$a$' times along y-axis.
    \item $f(x) \to \cfrac{1}{a} f(x); \quad a > 1$ : Stretch the graph of $f(x)$ `$a$' times along y-axis.
    
    \end{itemize}

        \includegraphics[width=0.9\linewidth,trim=0 0 0 0cm,clip]{figures/Signal_Stretch.png}

\end{frame}

\begin{frame}{Signal Stretching along $x$-axis}

    \begin{itemize}
    \item $f(x) \to  f(ax); \quad a > 1$ : Stretch the graph of $f(x)$ `$a$' times along x-axis.
    \item $f(x) \to f\bigg(\cfrac{1}{a} x\bigg); \quad a > 1$ : Shrink the graph of $f(x)$ `$a$' times along x-axis.
    
    \end{itemize}

        \includegraphics[width=0.9\linewidth,trim=0 0 0 0cm,clip]{figures/Signal_Sretch2.png}

\end{frame}

\begin{frame}{Example and MATLAB Code}


\begin{columns}
    \column{0.25\linewidth}
    $$
x(t) = 6t^2 + 4t + 70
$$
Code: \url{https://github.com/rahulbhadani/CPE381_FA26/blob/main/Code/Ch03_signal_operation.m}

        \column{0.75\linewidth}
        \includegraphics[width=0.99\linewidth,trim=0 0.0 0 1.0cm,clip]{figures/signal_operations.pdf}

\end{columns}
    
\end{frame}

\begin{frame}{Even and Odd Signals}
\begin{itemize}
    \item Even Signal: $x(t) = x(-t)$
\item Odd Signal:
$x(t) = -x(-t)$
\item Any signal can be represented by the sum of even and odd signals

$y(t) = y_e(t) + y_o(t)$

$$
y_e(t) = 0.5[ y(t)  + y(-t) ]$$

$$
y_o(t) = 0.5[y(t) - y(-t) ]
$$

\end{itemize}
\end{frame}

\begin{frame}{Periodic Signals}

\footnotesize

\begin{itemize}
    \item Defined for all possible values of $t, -\infty < t < \infty$.
    \item 
There is the real value $T_0 \in \Rbb^+$, called the fundamental frequency such that $x(  t + kT_0 ) = x(t), k \in \Ibb$.

\item A constant signal is periodic of a non-definable fundamental period.
\item A $\cos(\omega t + \theta)$, $\omega = 2\pi/T_0$ ,
$\omega = 2, \theta = -\pi/2, A= 2$.

\vspace{5pt}

\bf \color{red}
What’s the fundamental frequency, $1/T_0$?
\end{itemize}
\vfill

\end{frame}

\begin{frame}{Energy and Power of Signals}


    \begin{itemize}
        \item \textbf{Energy:}
        $$
        E_x = \int_{-\infty}^\infty | x(t)| ^2 dt
        $$

        \item \textbf{Power:}
        $$
        P_x = \lim_{T\to\infty}\cfrac{1}{2T} \int_{-T}^T | x(t)| ^2 dt
        $$
    \end{itemize}

    \bf A signal is called finite power if the signal power is finite.
\end{frame}


%%%%%%%%%%%%%%%%%%%%%%%%%%%%%%%%%%%%%%%%%%%%%%%%%%%%%
\subsection{Basic Signals as Building Blocks}

%-----------------------------------------

\begin{frame}
    \subsectionpage
\end{frame}

\begin{frame}{Complex Exponentials}
Consider $A = |A|e^{j\theta}$, $a = r + k\Omega_0$
\begin{itemize}
    \item $x(t) = Ae^{at}$ = ...
    \item Real part $f(t) = \Re{x(t)}, = ... $
    
    \vspace{20pt}
    
    $\quad  -|A|e^{rt} \leq f(t)) \leq |A|e^{rt}$. $r< 0$, $f(t)$ is damped, $r>0 $, $f(t)$ grows.
    \item  Imaginary part $g(t) = \Im{x(t)}, = ... $
    
    \vspace{20pt}
\end{itemize}
    
\end{frame}

\begin{frame}{Sinusoids}
A sinusoid of the general form:
$$
A\cos(\Omega_0 t + \theta) = A\sin (\Omega_0 t + \theta + \pi/ 2), \quad -\infty < t < \infty
$$

\begin{itemize}
    \item $A$ is the amplitude
    \item $\Omega_0 = 2\pi f_0$ is angular frequency in \texttt{rad/s}.
    \item $\theta$ is phase shift
    \item Fundamental period $T_0$ is 

    $$
    T_0 = \cfrac{2\pi}{\Omega_0} = \cfrac{1}{f_0}
    $$
\end{itemize}
    
\end{frame}

\begin{frame}{Rectangular pulse and Unit impulse}

\begin{columns}
    \column{0.7\linewidth}
    \begin{itemize}
        \item  A rectangular pulse of duration $\Delta$ and unit area:
        \begin{equation*}
            \begin{aligned}
                p_\Delta(t)= \begin{cases}
                    \cfrac{1}{\Delta},\quad -\Delta/2 \leq t \leq \Delta/2\\
                    0, \quad \textrm{otherwise}
                \end{cases}
            \end{aligned}
        \end{equation*}
        \item Unit Impulse:
        \begin{equation*}
            \delta(t) = \lim_{\Delta \to 0}p_\Delta (t)
        \end{equation*}
        \item \bf \color{red} Calculate 
        $$
        \int_{ -\infty}^{t}   p_\Delta(t)  dt
        $$
        Hint: there are three possible cases.
    \end{itemize}
    \column{0.3\linewidth}
\includegraphics[width=0.7\linewidth,trim=0 0 0 0cm,clip]{figures/unitpulse.png}
\end{columns}
\end{frame}

\begin{frame}{}
    
\end{frame}

\begin{frame}{Unit Step}

    \small
\begin{columns}
    \column{0.6\linewidth}
    \begin{itemize}
        \item  Integration of rectangular pulse:
        \begin{equation*}
            \begin{aligned}
                u_\Delta(t)= \int_{ -\infty}^{t}   p_\Delta(t) = \begin{cases}
                    1,\quad t \geq \cfrac{\Delta}{2}\\
                    \tfrac{1}{\Delta}(t + \tfrac{\Delta}{2}), \quad -\cfrac{\Delta}{2} \leq t \leq \cfrac{\Delta}{2}\\
                    0, \quad t < -\cfrac{\Delta}{2}
                \end{cases}
            \end{aligned}
        \end{equation*}
        \item Limit case:
        \begin{equation*}
            \lim_{\Delta \to 0}u_\Delta (t) = \begin{cases}
                    1,\quad t \geq 0\\
                    0, \quad t <0
                \end{cases}
        \end{equation*}

    \end{itemize}
    \column{0.4\linewidth}
    \begin{itemize}
            \item A common case is to ignore $t=0$ case, which gives us unit step function as
        \begin{equation*}
        u(t) = \begin{cases}
                    1,\quad t \geq 0\\
                    0, \quad t <0
                \end{cases}
        \end{equation*}
    \end{itemize}
\includegraphics[width=0.35\linewidth,trim=0 0 0 0cm,clip]{figures/unitstep.png}
\end{columns}
\end{frame}

\begin{frame}{Ramp Signal}
The ramp signal is $r(t) = t u(t)$
\begin{columns}
    \column{0.4\linewidth}
    \begin{itemize}
        \item The relation between the ramp, the unit step, and the unit impulse:
        \begin{equation*}
            \begin{aligned}
                \cfrac{dr(t)}{dt} & = u(t)\\
                \cfrac{d^2r(t)}{dt^2} & = \delta(t)
            \end{aligned}
        \end{equation*}
    \end{itemize}
    \column{0.6\linewidth}
    \includegraphics[width=0.65\linewidth,trim=0 0 0 0cm,clip]{figures/ramp_function.png}
    \end{columns}
\end{frame}

\begin{frame}{Triangular Pulse}
The triangular pulse is 
\begin{columns}
    \column{0.4\linewidth}
    \begin{equation*}
    \begin{aligned}
        \Lambda(t) = \begin{cases}
            t ,\quad 0 \leq t \leq 1\\
            -t + 2, \quad 1 < t \leq 2\\
            0,\quad \textrm{otherwise}
        \end{cases}
    \end{aligned}
    \end{equation*}
    \begin{itemize}
        \item $ \Lambda(t)$ can also be written as:
        \begin{equation*}
            \begin{aligned}
                \Lambda(t) =   r(t) - 2r(t-1) + r(t-2)
            \end{aligned}
        \end{equation*}
    \end{itemize}
        \column{0.6\linewidth}
        \includegraphics[width=0.85\linewidth,trim=0 0 0 0cm,clip]{figures/triangular_pulse.png}
    \end{columns}
\end{frame}

\begin{frame}[allowframebreaks]{Triangular Pulse as Ramp Functions}
Let's verify this by evaluating \( \Lambda(t) \) at different intervals:

\begin{columns}
    \column{0.5\linewidth}
        First, find out the first part:

    \begin{equation*}
            \begin{aligned}
                \Lambda_1(t) =  \begin{cases}
            t ,\quad 0 \leq t \leq 1\\
            0,\quad \textrm{otherwise}
        \end{cases}
            \end{aligned}
        \end{equation*}
    which can be written as:

can be written using the ramp function \( r(t) \) as:
    \begin{equation*}
        \Lambda_1(t) = r(t) - r(t-1)
    \end{equation*}
    \column{0.5\linewidth}
    \includegraphics[width=0.75\linewidth,trim=0 0 0 0cm,clip]{figures/lambda_1.png}
\end{columns}

\begin{columns}
    \column{0.5\linewidth}
    Second part:
        \begin{equation*}
            \begin{aligned}
                \Lambda_1(t) =  \begin{cases}
            -t + 2, \quad 1 < t \leq 2\\
            0,\quad \textrm{otherwise}
        \end{cases}
            \end{aligned}
        \end{equation*}

can be written using the ramp function \( r(t) \) as:
    \begin{equation*}
        \Lambda_1(t) = r(t-2) - r(t-1)
    \end{equation*}
    \column{0.5\linewidth}
    \includegraphics[width=0.95\linewidth,trim=0 0 0 0cm,clip]{figures/lambda_2.png}
\end{columns}

    Putting together
    \includegraphics[width=0.40\linewidth,trim=0 0 0 0cm,clip]{figures/lambda_1_2.png}

    \bigskip

    \small
    
    Code for the graph: \url{https://github.com/rahulbhadani/CPE381_FA26/blob/main/Code/Ch03_ramp.m}

\end{frame}

\begin{frame}{Sifting Property}
The product of $f (t)$ and $\delta(t)$ gives zero everywhere except at the origin where we get an impulse of area $f (0)$, that is, $f (t)\delta(t) = f (0)\delta(t)$.

Hence,

$$
\int_{-\infty}^\infty f(t)\delta(t)dt = \int_{-\infty}^\infty f(0)\delta dt = f(0) \int_{-\infty}^\infty\delta (t) dt  = f(0)
$$
This is called \textbf{Sifting Property}.

If we delay or advance the $\delta(t)$ function in the integrated, the result is that all values of $f (t)$ are sifted out except for the value corresponding to the location of the delta function, that is,

$$
\int_{-\infty}^\infty f(t)\delta(t - \tau)dt  = \int_{-\infty}^\infty f(\tau)\delta(t- \tau) dt  = f(\tau) \int_{-\infty}^\infty \delta(t-\tau) dt  = f(\tau)\quad \text{for any }\tau
$$

\end{frame}

\begin{frame}{Generic Representation of Signals}
Hence, if we do integration in terms of variable $\tau$, we get a generic representation of signals in terms of impulse and shifted impulse.
\begin{columns}
    \column{0.45\linewidth}
    $x(t) = \int_{-\infty}^\infty x(\tau ) \delta(t-\tau)d\tau $

        \includegraphics[width=0.85\linewidth,trim=0 0 0 0cm,clip]{figures/signal_multiplication.png}

    \column{0.55\linewidth}
    \includegraphics[width=0.85\linewidth,trim=0 0 0 0cm,clip]{figures/signal_sums.png}
        \begin{itemize}
            \item Approximation of $x(t)$:
            \footnotesize
            {
            \begin{equation*}
                x_\Delta(t) = \sum_{-\infty}^\infty x_\Delta(t-k\Delta) = \sum_{-\infty}^\infty x(k\Delta) p_\Delta(t-k\Delta)\Delta
            \end{equation*}
            }
        \end{itemize}
\end{columns}
In the limit as $\Delta \to 0$ these pulses become impulses, separated by an infinitesimal value:

$\lim_{\Delta \to 0 }x_\Delta(t) \to x(t)  =  \int_{-\infty}^\infty x(\tau ) \delta(t-\tau)d\tau $
    
\end{frame}

%%%%%%%%%%%%%%%%%%%%%%%%%%%%%%%%%%%%%%%%%%%%%%%%%%%%%
\subsection{Modulation and Windowing}

%-----------------------------------------

\begin{frame}
    \subsectionpage
\end{frame}

\begin{frame}{Modulation}
    Multiplication by a complex exponential shifts the frequency of the original signal.
    \begin{block}{Definition}
        Superimposing a low-frequency signal on a high-frequency carrier signal is called \textbf{Modulation}.
    \end{block}
    
    \textbf{Example:}

    Consider an exponential signal $x(t) = e^{j\Omega_0 t}$ of frequency $\Omega_0$. If we multiply an exponential $e^{j\phi t}$ with $x(t)$, then:
    $$
    x(t)e^{j\phi t} = e^{j(\Omega_0 + \phi) t} = \cos((\Omega_0 + \phi) t) + j \sin((\Omega_0 + \phi) t)
    $$

    $
    \phi > 0: 
    $ the frequency of new exponential is greater than $\Omega_0$, otherwise lower.
    
\end{frame}

\begin{frame}{Various Types of Modulation}
$A(t)cos( \Omega(t)t + \theta(t) )$
    \begin{itemize}
        \item $A(t)$ changes: Amplitude Modulation
\item $\Omega(t)$ changes: Frequency Modulation
\item $\theta(t)$ changes: Phase Modulation
    \end{itemize}
\end{frame}

\begin{frame}{Windowing}

\begin{columns}
    \column{0.5\linewidth}
        For a window signal $w(t)$, the time-windowed signal 
        
        $x(t)w(t)$ displays $x(t)$ within the support of $w(t)$.
        
    \textbf{Example:}

    $x(t) = \sin(2\pi t)$

    \begin{equation*}
        w(t) = \begin{cases} 1 \quad \text{if } -0.5 \leq t \leq 0.5 \\ 0 \quad \text{otherwise} \end{cases}
    \end{equation*}

\small

    Code for the graph: \url{https://github.com/rahulbhadani/CPE381_FA26/blob/main/Code/Ch03_windowing.m}
    \column{0.5\linewidth}
            \includegraphics[width=0.65\linewidth,trim=0 0 0 0cm,clip]{figures/signal_windowing.pdf}
\end{columns}

\end{frame}

\begin{frame}{Classwork}
    
\end{frame}

\begin{frame}{Classwork}
    
\end{frame}
\begin{frame}{Classwork}
    
\end{frame}
\begin{frame}{Classwork}
    
\end{frame}
\begin{frame}{Classwork}
    
\end{frame}
\begin{frame}{Classwork}
    
\end{frame}
\begin{frame}{Classwork}
    
\end{frame}
\begin{frame}{Classwork}
    
\end{frame}
\begin{frame}{Classwork}
    
\end{frame}


\subsection{Linearity of Signals}

\begin{frame}
    \subsectionpage
\end{frame}

\begin{frame}{Properties of Linearity}

    A signal $f$ is linear if it has these properties,

    \begin{enumerate}
        \item Homogeneity (scaling): $f(at) = af(t)$ where $a$ is a constant scalar.
        \item Additivity: $f(t_1+t_2) = f(t_1) + f(t_2)$.
    \end{enumerate}
    
    \bigskip

    \centering

    The homogeneity and additivity together are called the \textbf{principle of superposition}
\end{frame}

\begin{frame}{Example 1}

    Consider a signal $f(t) = 5t$.  Does it satisfy homogeneity?

    \bigskip

    \textbf{Solution:}

    \begin{enumerate}
        \item Apply the function to a scaled input, i.e. $a \cdot t$: $$f(a\cdot t) = 5 · (a\cdot t) = 5a t$$
        \item Now, scale the output $f(t)$ by the same factor $a$: $$a\cdot f(t) = a\cdot  (5t) = 5at$$
        \item LHS = RHS. Hence it satisfies homogeneity.
    \end{enumerate}
    
\end{frame}

\begin{frame}{Example 2}
    
    Consider a signal $f(t) = 5t$.  Does it satisfy additivity?

    \bigskip

    \textbf{Solution:}

    \begin{enumerate}
        \item Apply the function to the sum of two inputs $t_1$ and $t_2$: $$f(t_1 + t_2) = 5 · (t_1 + t_2) = 5t_1 + 5t_2$$
        \item Now, add the outputs of the function applied to each input individually: $$ f(t_1) + f(t_2) = (5t_1) + (5t_2) = 5t_1 + 5t_2 $$
        \item LHS = RHS. Hence it satisfies additivity.
    \end{enumerate}


\end{frame}

\begin{frame}{Example 3}
    

    Consider a signal $f(t) = 5t+3$.  Does it satisfy homogeneity?

    \vfill

\end{frame}

\begin{frame}{}
    
\end{frame}

\begin{frame}{Example 4}
    

    Consider a signal $f(t) = t^2$.  Does it satisfy additivity?

    \vfill

\end{frame}

\begin{frame}{}
    
\end{frame}


\section{Continuous-Time Systems}

\begin{frame}
    \sectionpage
\end{frame}


\begin{frame}{System Thinking}
    \begin{center}
        
    \textbf{System is an abstraction.}
    
    Think of a device or a process as a system -- something that transforms one signal into another.
    
    \centering
      \includegraphics[width=0.65\linewidth,trim=0 0 0 0cm,clip]{figures/system.pdf}

      \end{center}
      
\end{frame}

\begin{frame}{Continuous-time Systems}
    \begin{center}
        \bf \Large
        When inputs and outputs are continuous, the system is continuous.
    \end{center}
\end{frame}

%%%%%%%%%%%%%%%%%%%%%%%%%%%%%%%%%%%%%%%%%%%%%%%%%%%%%
\subsection{Linear Time-Invariant (LTI) Continuous-Time Systems}

%-----------------------------------------

\begin{frame}
    \subsectionpage
\end{frame}

\begin{frame}{Linear Time Invariant (LTI) Continuous-Time Systems}
    For most analyses at system-level, we care about the following properties:
    \begin{itemize}
        \item Linearity
        \item Time-invariance
        \item Causality
        \item Stability
    \end{itemize}
\end{frame}

\begin{frame}{Linearity of a System}
   \large
   \centering
   \begin{equation*}
\mathcal{S}[a \cdot x(t) + b \cdot y(t)] = \mathcal{S}[a \cdot x(t)] + \mathcal{S}[b \cdot y(t)] = a \cdot \mathcal{S}[x(t)] + b \cdot \mathcal{S}[y(t)]
\end{equation*}
\end{frame}

\begin{frame}{Time-invariance of a System}
     \large
   If
   \begin{equation*}
x(t) \Rightarrow y(t) = \mathcal{S}[x(t)],
\end{equation*}

then, the  time-invariance says:
\begin{equation*}
x(t \mp \tau) \Rightarrow y(t \mp \tau) = \mathcal{S}[x(t \mp \tau)].
\end{equation*}
\end{frame}

\begin{frame}{Linear Time Invariant (LTI) Systems}
    \begin{center}
        \bf \large 
        \begin{gradbox}{}
            A system that satisfies both linearity and time-invariance is an LTI system.
        \end{gradbox}
    \end{center}
\end{frame}

\begin{frame}{Class Problem 1}
    Consider a system 
    $g(t) = \mathcal{S}[f(t)] = a [f(t)]^2 - b [f(t)]$
    \begin{itemize}
        \item Is this system linear?
        \item Is this system time-invariant?
    \end{itemize}
\end{frame}

\begin{frame}{Class Problem 2}
    Consider 
    \begin{equation*}
        g(y) = \Scal[f(t)] = \int_{-\infty}^\infty  f(t)\textrm{rect}(y-t)dt
    \end{equation*}
    \begin{itemize}
        \item Is this system linear?
        \item Is this system time-invariant?
    \end{itemize}
\end{frame}

\begin{frame}{Static Systems vs Dynamic Systems}
    \begin{gradbox}{}
        In static systems, output depends on the current input but not on the historical input or output
    \end{gradbox}
\vspace{10pt}
    \begin{gradblock}{}
    In dynamic systems, output depends on current input, past inputs, and past outputs.
    \end{gradblock}
\vspace{10pt}
       \begin{gradbox}{}
       Differential equations are used to represent dynamic systems (literature also refers to them as dynamical systems).
    \end{gradbox}

\end{frame}

\begin{frame}{Representation of Dynamic Systems}
Dynamic systems are represented by ordinary differential equations (ODE) with constant coefficients:
\small
\begin{equation*}
\begin{aligned}
    a_n \frac{d^n y(t)}{dt^n} + & a_{n-1} \frac{d^{n-1} y(t)}{dt^{n-1}} + \cdots + a_1 \frac{dy(t)}{dt} + a_0 y(t) = \\
   &  b_m \frac{d^m x(t)}{dt^m} + b_{m-1} \frac{d^{m-1} x(t)}{dt^{m-1}} + \cdots + b_1 \frac{dx(t)}{dt} + b_0 x(t)
\end{aligned}
\end{equation*}
with $n$ initial conditions, $x(t) = 0$ for $t < 0$.

\bf
\begin{center}
    This equation represents a linear ordinary differential equation with constant coefficients, where $ y(t)$  is the output,  $x(t)$  is the input, and  $a_i$  and  $b_i$  are the constant coefficients. 
\end{center}
\end{frame}

\begin{frame}{Response of a Dynamic Systems}
    Assessing the output of the system:
    \begin{itemize}
        \item Consider output due to input only, considering initial conditions to be zero.  \textbf{(zero-state response)}.
        \item Consider output due to initial conditions only, considering inputs to be zero. \textbf{(zero-input response)}.
        \item Complete response is the sum of the above two.


    \end{itemize}
\end{frame}

\begin{frame}{Differential Equations for Ideal Systems}
\begin{columns}
    \column{0.33\linewidth}
    Electrical Inductance
        \centering
      \includegraphics[width=0.8\linewidth,trim=0 0 0 0cm,clip]{figures/electrical_inductance.png}
 \column{0.33\linewidth}
 Electrical Capacitance
        \centering
      \includegraphics[width=0.8\linewidth,trim=0 0 0 0cm,clip]{figures/electrical_capacitance.png}
 \column{0.33\linewidth}
 Electrical Resistance
        \centering
      \includegraphics[width=0.8\linewidth,trim=0 0 0 0cm,clip]{figures/electrical_resistance.png}
      
\end{columns}
    
\end{frame}

\begin{frame}{Parallel RLC Circuit}
\begin{columns}
    \column{0.5\linewidth}
    $$
    \cfrac{v(t)}{R} + C \cfrac{dv(t)}{dt} + \cfrac{1}{L}\int_0^t v(t)dt = r(t)
    $$
    \column{0.5\linewidth}
     \includegraphics[width=0.99\linewidth,trim=0 0 0 0cm,clip]{figures/parallel_RLC.png}
      
\end{columns}
    
\end{frame}

\begin{frame}{Damped Mass-spring Oscillator}
\begin{columns}
    \column{0.5\linewidth}
    $k = $ stiffness of the spring

    $b = $ damping coefficient

    $m = $ mass of the block

    $y = y(t)$ time-dependent displacement from Equilibrium point.

    $
    my'' = F - ky - by'
    $

    \column{0.5\linewidth}
 \includegraphics[width=0.75\linewidth,trim=0 0 0 0cm,clip]{figures/mass_spring_oscillator.png}
    
\end{columns}
    
\end{frame}

\subsection{Convolutional Integral}

%-----------------------------------------

\begin{frame}
    \subsectionpage
\end{frame}

\begin{frame}{Response to an Impulse Signal}

\large
\begin{center}
    Output when a system is presented with a unit-impulse signal is called \textbf{response to an impulse signal}, or in short \textbf{impulse response}.
\end{center}
    
\end{frame}

\begin{frame}{Why is Impulse Response Useful?}
It allows us to predict what the system's output will look like in the \textbf{time-domain}.

\begin{gradblock}{}

Do you remember generalized signal representation using impulse signals?
\LARGE
$$
x(t) = \int_{-\infty}^\infty x(\tau) \delta (t - \tau)d \tau
$$
  
\end{gradblock}

    \color{red}\bf What would we get if we pass $x(t)$ through an LTI system $\Scal$?
\end{frame}

\begin{frame}{Impulse Response}
What would we get if we pass $x(t)$ through an LTI system $\Scal$?

\vspace{10pt}

We get

\begin{gradyellow}
\Large
    $$
    y(t) = \int_{-\infty}^\infty x(\tau) h (t - \tau)d \tau
    $$
\end{gradyellow}

\begin{columns}
    \column{0.4\linewidth}
    where $h(t)$ is the impulse response. Out is the superposition of the responses of each term in the figure.
    \column{0.6\linewidth}
    \includegraphics[width=0.88\linewidth,trim=0 0 0 0cm,clip]{figures/signal_sums.png}
\end{columns}




\end{frame}

\begin{frame}{Definition of Convolution Integral}

The response of an LTI system $\Scal$, represented by its impulse response $h(t)$, to any signal $x(t)$ is the \textbf{convolution integral}.

\begin{gradgreen}
    
    $$
    y(t) =  \int_{-\infty}^\infty  x(\tau) h(t-\tau)d\tau =  \int_{-\infty}^\infty  x(t-\tau) h(\tau) d\tau  = [x * h](t) = [h*x](t)
    $$
    
\end{gradgreen}
    where $*$ stands for the convolution integral of the input signal $x(t)$ and the impulse response $h(t)$ of the system.
\end{frame}

\begin{frame}{Some Points on Convolution Integral}

\Large
\begin{itemize}
    \item Impulse response is fundamental to characterizing an LTI system.
    \item A system represented by the convolution integral is LTI. It is a general representation of an LTI system.
\end{itemize}
\end{frame}
\begin{frame}{Putting Multiple Systems Together}
\begin{columns}
    \column{0.5\linewidth}
    \includegraphics[width=0.7\linewidth,trim=0 0 0 0cm,clip]{figures/connected_systems.png}
    \column{0.5\linewidth}

\Large 
    A: $h(t) = [h_1 * h_2] (t) = [h_2 * h_1](t)$

    B: $h(t) = h_1(t) + h_2(t)$

    C: $h(t) = [h_1 - h*h_1*h_2](t)$
\end{columns}    
\end{frame}

%%%%%%%%%%%%%%%%%%%%%%%%%%%%%%%%%%%%%%%%%%%%%%%%%%%%%
\subsection{Causality and Stability}

%-----------------------------------------

\begin{frame}
    \subsectionpage
\end{frame}

\begin{frame}{Causality}
    \begin{itemize}
        \item Cause-effect relationship,
        \item A system is causal if
        \begin{gradred}
            \begin{itemize}
                \item whenever its input $x(t) = 0$, and there is no initial condition, $y(t) = 0$.
                \item the output $y(t)$ doesn't depend on the future inputs.
            \end{itemize}
        \end{gradred}
        \item Causality is independent of linearity and time-invariance.
    \end{itemize}
    \bf LTI system with impulse response $h(t)$ is causal if $h(t) = 0$ for $t < 0$, and 
    \begin{gradblue}
        $$
    y(t) = \int_0^t x(\tau) h(t-\tau)d\tau
    $$
    \end{gradblue}
    
\end{frame}

\begin{frame}{Bounded Input Bounded Output Stability}
\begin{itemize}
    \item If input is bounded, i.e. $ |x(t)| \leq M, M  < \infty$, then output is bounded under certain conditions.
    \item Condition for bounded output is
\begin{gradyellow}
   $$ |y(t)| = \bigg| \int_{-\infty}^\infty x(t-\tau)h(\tau)d\tau \bigg| \leq \int_{-\infty}^\infty|x(t-\tau)| | h(\tau)| d\tau \leq M \int_{-\infty}^\infty |h(\tau)| d\tau \leq MK <\infty $$

   where the integral $\int_{-\infty}^\infty |h(\tau)|d\tau \leq K < \infty$

i.e., the impulse is absolutely integrable.
\end{gradyellow}
\end{itemize}
\end{frame}

\begin{frame}{To Calculate if an LTI system is BIBO Stable}

\Large
    \begin{itemize}
        \item Compute its impulse response $h(t)$.
        \item Integrate it over $-\infty$ to $\infty$, effectively over the given support of the signal.
        \item If the integration is finite, the system is BIBO Stable.
    \end{itemize}
\end{frame}

\begin{frame}{Classwork}
    
\end{frame}

\begin{frame}{Classwork}
    
\end{frame}
\begin{frame}{Classwork}
    
\end{frame}
\begin{frame}{Classwork}
    
\end{frame}
\begin{frame}{Classwork}
    
\end{frame}
\begin{frame}{Classwork}
    
\end{frame}
\begin{frame}{Classwork}
    
\end{frame}
\begin{frame}{Classwork}
    
\end{frame}
\begin{frame}{Classwork}
    
\end{frame}
\begin{frame}{Classwork}
    
\end{frame}


\begin{frame}{Up Next}
\begin{itemize}
	\item Fourier Series
%    \item Continuous-time Systems
%    \begin{itemize}
%        \item Linear-Time Invariance
%        \item Static vs Dynamic Systems
%        \item Convolutional Integral
%        \item BIBO Stability
%    \end{itemize}
\end{itemize}
    
\end{frame}
\end{document}